\begin{table}[H]
\centering
\caption{PIB por actividad económica, 25 agrupaciones CIIU, 2010--2025}
\label{tab:valor_agregado_25_agrupaciones}
\small
\begin{tabular}{p{0.50\textwidth}rrr}
\toprule
Actividad económica & 2010 & 2025 & Crec. \\
\midrule
Arte, entretenimiento, recreación y otros servicios & 11,0 & 40,6 & 9,09\% \\
Financieras y seguros & 22,3 & 53,1 & 5,95\% \\
Salud y servicios sociales & 23,6 & 49,8 & 5,11\% \\
Administración pública & 31,3 & 62,7 & 4,75\% \\
Comercio y reparación & 51,2 & 89,0 & 3,76\% \\
Información y comunicaciones & 18,3 & 31,3 & 3,66\% \\
Alojamiento y comida & 23,8 & 40,6 & 3,63\% \\
Educación & 31,1 & 52,7 & 3,59\% \\
Agropecuario & 39,9 & 63,9 & 3,18\% \\
Hogares como empleadores & 4,3 & 6,7 & 3,06\% \\
Profesionales y administrativas & 45,4 & 70,1 & 2,94\% \\
Transporte y almacenamiento & 32,8 & 49,8 & 2,81\% \\
Inmobiliarias & 59,9 & 90,1 & 2,75\% \\
Electricidad, gas y vapor & 14,9 & 21,9 & 2,60\% \\
Alimentos, bebidas y tabaco & 23,9 & 33,5 & 2,26\% \\
Refinación, químicos y minerales & 30,8 & 39,6 & 1,69\% \\
Textiles, confecciones y cuero & 9,6 & 11,9 & 1,42\% \\
Obras civiles & 9,8 & 11,9 & 1,36\% \\
Muebles y otras manufactureras & 5,0 & 6,0 & 1,28\% \\
Agua, saneamiento y desechos & 7,0 & 8,4 & 1,19\% \\
Madera, papel e impresión & 5,4 & 6,4 & 1,17\% \\
Metalurgia, maquinaria y equipo & 13,1 & 15,5 & 1,10\% \\
Actividades especializadas de construcción & 8,5 & 9,1 & 0,41\% \\
Edificaciones & 22,0 & 20,5 & -0,49\% \\
Minas & 38,4 & 34,7 & -0,67\% \\
\bottomrule
\end{tabular}
\caption*{\footnotesize Nota: valores en billones de pesos constantes de 2015. El cuadro usa las 25 agrupaciones de actividad económica del DANE por el enfoque de producción. En sentido estricto, el numerador corresponde al valor agregado bruto de cada actividad. Fuente: cálculos propios con DANE, anexo de producción a precios constantes.}
\end{table}

\textbf{La desagregación a 25 agrupaciones muestra que buena parte de la heterogeneidad está dentro de las grandes actividades económicas.} Las mayores tasas de crecimiento del PIB se observan en Arte, entretenimiento, recreación y otros servicios (9,09\%); Financieras y seguros (5,95\%); Salud y servicios sociales (5,11\%). En el otro extremo, las menores tasas se registran en Minas (-0,67\%); Edificaciones (-0,49\%); Actividades especializadas de construcción (0,41\%). Esta lectura ayuda a identificar dónde se originan los cambios agregados: no basta con saber que una gran actividad crece o cae, porque dentro de ella pueden coexistir subactividades con trayectorias muy distintas.
